\documentclass[a4paper,12pt]{article}
% Paquetes necesarios
\usepackage[utf8]{inputenc}  % Para caracteres en español
\usepackage{amsmath, amssymb} % Para símbolos matemáticos
\usepackage{fancyhdr} % Para encabezado y pie de página
\usepackage{geometry} % Para márgenes
\usepackage{enumitem} % Para modificar listas
\usepackage{tabularx}
\usepackage{calc}
\usepackage{array}

\newcolumntype{L}[1]{>{\raggedright\arraybackslash}p{#1}} % Columna izquierda, alineada a la izquierda
\newcolumntype{R}[1]{>{\raggedleft\arraybackslash}p{#1}}  % Columna derecha, alineada a la derecha

% Configuración de márgenes
\geometry{left=2.5cm, right=2.5cm, top=2.5cm, bottom=3cm}

% Configuración del encabezado y pie de página
\pagestyle{fancy}
\lfoot{Ing. Luis A. Molina}  % Pie de página a la izquierda
\cfoot{MAT207}  % Centro vacío
\rfoot{\thepage} % Número de página a la derecha

% Inicio del documento
\begin{document}

% Título de la práctica
\begin{center}
    \Large\textbf{Práctica 05 - Ecuaciones Diferenciales Ordinarias de Orden Superior: Variación de Parámetros}\\[1cm]  
\end{center}
  \textbf{Nombre:} \rule{9cm}{0.4pt}  \textbf{Fecha:} \rule{4.5cm}{0.4pt}

\vspace{0.5cm}

\subsubsection*{A. Resuelva las siguientes ecuaciones diferenciales utilizando el método de variación de parámetros}

\begin{enumerate}
    % Ejercicios Fáciles
    \item
    \begin{tabularx}{\linewidth}{L{0.45\linewidth} R{0.45\linewidth}}
    \( y'' + y = \sec x \) & \(\displaystyle y = c_1 \cos x + c_2 \sin x + \cos x \ln|\cos x| + x \sin x \)
    \end{tabularx}
    
    \item
    \begin{tabularx}{\linewidth}{L{0.45\linewidth} R{0.45\linewidth}}
    \( y'' - 2y' + y = \dfrac{e^x}{x} \) & \(\displaystyle y = c_1 e^x + c_2 x e^x + x e^x \ln|x| \)
    \end{tabularx}
    
    % Ejercicios Medios
    \item
    \begin{tabularx}{\linewidth}{L{0.45\linewidth} R{0.45\linewidth}}
    \( y'' + y = \tan x \) & \(\displaystyle y = c_1 \cos x + c_2 \sin x - \cos x \ln|\sec x + \tan x| \)
    \end{tabularx}
    
    \item
    \begin{tabularx}{\linewidth}{L{0.45\linewidth} R{0.45\linewidth}}
    \( y'' + 4y = \csc 2x \) & \(\displaystyle y = c_1 \cos 2x + c_2 \sin 2x - \frac{1}{2} x \cos 2x + \frac{1}{4} \sin 2x \ln|\sin 2x| \)
    \end{tabularx}
    
    % Ejercicios Difíciles
    \item
    \begin{tabularx}{\linewidth}{L{0.45\linewidth} R{0.45\linewidth}}
    \( y'' + y = \sec^3 x \) & \(\displaystyle y = c_1 \cos x + c_2 \sin x + \frac{1}{2} \sec x \)
    \end{tabularx}
    
    \item
    \begin{tabularx}{\linewidth}{L{0.45\linewidth} R{0.45\linewidth}}
    \( y'' - 2y' + 2y = e^x \sec x \) & \(\displaystyle y = e^x(c_1 \cos x + c_2 \sin x) + e^x \cos x \ln|\cos x| + x e^x \sin x \)
    \end{tabularx}

    \item
    \begin{tabularx}{\linewidth}{L{0.45\linewidth} R{0.45\linewidth}}
    \( x^2 y'' - 2x y' + 2y = x^3 \ln x \) & \(\displaystyle y = c_1 x + c_2 x^2 + \frac{1}{2} x^3 \ln x - \frac{3}{4} x^3 \)
    \end{tabularx}
    
\end{enumerate}

\end{document}
